\documentclass[12pt,a4paper]{article}
\usepackage[utf8]{inputenc}
\usepackage[T5]{fontenc}
\usepackage[vietnamese]{babel}
\usepackage[margin=2.5cm]{geometry}
\usepackage{graphicx}
\usepackage{longtable}
\usepackage{booktabs}
\usepackage{listings}
\usepackage{xcolor}
\usepackage{hyperref}
\usepackage{enumitem}

\lstset{
  basicstyle=\ttfamily\small,
  breaklines=true,
  columns=fullflexible,
  frame=single,
  backgroundcolor=\color{gray!5},
}

\title{\textbf{BÁO CÁO ĐỒ ÁN GIỮA KỲ\\
Xây dựng ngữ liệu đơn ngữ chữ Quốc ngữ\\
chuyên ngành lịch sử Việt Nam (Track A -- HVQ)}}
\author{Mã tác phẩm: \texttt{HVQ\_036}}
\date{\today}

\begin{document}
\maketitle
\tableofcontents

\section{Giới thiệu}

Đồ án thuộc hạng mục \textbf{A. HVQ -- Xây dựng ngữ liệu đơn ngữ chữ Quốc ngữ
chuyên ngành lịch sử Việt Nam}. Đầu vào là văn bản (text-based PDF, không phải
ảnh scan), do đó theo yêu cầu đề bài chỉ cần thực hiện hai bước: \emph{tách câu}
và \emph{gán nhãn thực thể (NER)}; bước hiệu đính chính tả (OCR correction)
được bổ sung thêm bằng LLM để nâng cao chất lượng ngữ liệu do văn bản gốc có
nhiều lỗi chính tả/dấu câu từ quá trình số hóa.

Ngữ liệu đầu vào gồm 2 quyển sách lịch sử Việt Nam, được gán chung mã tác phẩm
\texttt{HVQ\_036} và phân biệt theo mã chương/quyển:

\begin{itemize}
  \item \textbf{HVQ\_036\_LLKN}: \emph{Lịch Sử Việt Nam Từ Lê Lợi Khởi Nghĩa
    Đến Nguyễn Suy Vong} -- 165 trang PDF (158 trang nội dung sau khi loại bỏ
    mục lục/bìa theo bookmark).
  \item \textbf{HVQ\_036\_TTLS}: \emph{Tiến Trình Lịch Sử Việt Nam} -- 396
    trang PDF, không có bookmark nên lấy toàn bộ nội dung.
\end{itemize}

\section{Yêu cầu bài toán}

Theo \texttt{OutputRequirement.pdf} (mục A -- HVQ), với đầu vào dạng text, mỗi
quyển/tập cần sinh ra:

\begin{itemize}
  \item \texttt{[matacpham]\_seg.tsv}: nội dung đã tách câu, định dạng
    \texttt{[sentence\_id]\textbackslash t[sentence]}.
  \item \texttt{[matacpham]\_ner.json}: kết quả NER, mỗi câu là một object
    \texttt{\{sentence\_id, sentence, entities\}}, bộ nhãn tối thiểu gồm
    \texttt{PER, LOC, ORG, TITLE, TME, NUM} (ví dụ trong đề còn dùng thêm nhãn
    \texttt{DYNASTY}).
\end{itemize}

Nếu ngữ liệu gồm nhiều quyển/tập, mỗi quyển được lưu trong một thư mục riêng,
tên file mang tiền tố \texttt{[matacpham\_chapter]}:

\begin{lstlisting}
HVQ_036/
 |-- HVQ_036_LLKN/
 |    |-- HVQ_036_LLKN_seg.tsv
 |    |-- HVQ_036_LLKN_ner.json
 |    `-- HVQ_036_LLKN_corrections.jsonl
 `-- HVQ_036_TTLS/
      |-- HVQ_036_TTLS_seg.tsv
      |-- HVQ_036_TTLS_ner.json
      `-- HVQ_036_TTLS_corrections.jsonl
\end{lstlisting}

\section{Công cụ sử dụng}

\begin{itemize}
  \item \textbf{Trích xuất text}: PyMuPDF (\texttt{fitz}) -- đọc bookmark cấp 1
    của PDF để xác định phạm vi trang nội dung chính (bỏ mục lục, lời nói đầu,
    bìa), sau đó trích text thô theo từng trang.
  \item \textbf{Hiệu đính OCR}: Google Gemini (\texttt{gemini-flash-lite-latest},
    qua SDK \texttt{google-genai}) -- sửa lỗi chính tả/dấu câu và khôi phục dấu
    câu cuối câu trước khi tách câu.
  \item \textbf{Tách câu}: \texttt{underthesea} (\texttt{sent\_tokenize}), kết
    hợp ranh giới khối văn bản (đầu quyển/chương) làm ranh giới câu cứng.
  \item \textbf{NER}: kết hợp mô hình CRF chung của \texttt{underthesea} (cho
    nhãn \texttt{PER/LOC/ORG}) với một tầng luật + gazetteer tự xây dựng, chuyên
    biệt cho văn bản lịch sử Việt Nam (nhãn \texttt{TME}, \texttt{DYNASTY},
    \texttt{TITLE}, \texttt{NUM}, và bổ sung \texttt{LOC/PER} mà mô hình chung
    bỏ sót).
\end{itemize}

\section{Kiến trúc pipeline}

Pipeline gồm 4 giai đoạn chính, cộng một giai đoạn hiệu đính tùy chọn xen giữa
bước 2 và 3:

\begin{enumerate}
  \item \textbf{Trích xuất} (\texttt{pipeline/extract.py}) --
    \texttt{detect\_content\_range()} đọc bookmark cấp 1 để xác định trang bắt
    đầu/kết thúc nội dung chính; \texttt{extract\_pages()} trích text theo
    trang (đánh số từ 1).
  \item \textbf{Chuẩn hóa} (\texttt{pipeline/normalize.py}) --
    \texttt{normalize\_page\_text()}: chuẩn hóa Unicode (NFC), chuyển dấu thanh
    kiểu cũ sang kiểu mới (\texttt{hoà}$\to$\texttt{hòa}, \texttt{thuỷ}$\to$
    \texttt{thủy}, có bảo vệ cho trường hợp \texttt{qu}), sửa lỗi đặt sai dấu
    do OCR (\texttt{tòan}$\to$\texttt{toàn}), loại bỏ header lặp lại/số
    trang/URL/tiêu đề in hoa/nhiễu OCR, và nối lại các dòng bị ngắt do xuống
    dòng hoặc gạch nối. Tiêu đề quyển/phần được loại khỏi nội dung nhưng vẫn
    giữ vai trò ranh giới khối.
  \item \textbf{Hiệu đính OCR (tùy chọn, \texttt{--correct})}
    (\texttt{pipeline/correct.py}) -- chia mỗi trang đã chuẩn hóa thành các
    đoạn $\le 2500$ ký tự, gửi từng đoạn cho Gemini để sửa lỗi chính tả/dấu
    câu. Một ngưỡng bảo vệ theo tỉ lệ độ dài (\emph{length-ratio guard}, mặc
    định 30\%) sẽ từ chối kết quả lệch quá xa so với bản gốc, giữ nguyên đoạn
    gốc trong trường hợp đó. Kết quả được cache trên đĩa theo khóa
    \texttt{hash(model, đoạn văn)} để lần chạy sau không gọi lại API cho các
    đoạn đã xử lý. Nếu gọi API lỗi (hết quota, lỗi mạng...), pipeline lùi về
    dùng đoạn gốc (\texttt{fallback\_error}) thay vì dừng chương trình. Toàn bộ
    cặp gốc/đã sửa và hành động tương ứng được ghi vào
    \texttt{<code>\_corrections.jsonl}.
  \item \textbf{Tách câu} (\texttt{pipeline/segment.py}) --
    \texttt{split\_sentences()} dùng \texttt{underthesea.sent\_tokenize} trên
    từng khối văn bản đã chuẩn hóa, coi ranh giới khối là ranh giới câu cứng,
    loại bỏ các câu quá ngắn hoặc không chứa chữ cái.
  \item \textbf{NER} (\texttt{pipeline/ner.py}) --
    \texttt{extract\_entities()} hợp nhất hai nguồn: mô hình CRF chung của
    underthesea và tầng luật/gazetteer lịch sử. Các span chồng lấn được xử lý
    theo nguyên tắc span dài hơn thắng, sau đó theo thứ tự ưu tiên cố định:
    \texttt{TME > DYNASTY > TITLE > model > rule-LOC > rule-PER > NUM}.
\end{enumerate}

\section{Kết quả}

\subsection{Số liệu tổng quan}

\begin{table}[h]
\centering
\begin{tabular}{lrrrr}
\toprule
Quyển & Số trang & Số câu & Số thực thể & Số đoạn hiệu đính \\
\midrule
HVQ\_036\_LLKN & 158 & 1{,}498 & 4{,}472 & 153 \\
HVQ\_036\_TTLS & 396 & 7{,}276 & 16{,}802 & 394 \\
\midrule
\textbf{Tổng} & 554 & 8{,}774 & 21{,}274 & 547 \\
\bottomrule
\end{tabular}
\caption{Số liệu tổng quan hai quyển sau khi xử lý toàn bộ pipeline.}
\end{table}

\subsection{Phân bố nhãn thực thể}

\begin{table}[h]
\centering
\begin{tabular}{lrr}
\toprule
Nhãn & HVQ\_036\_LLKN & HVQ\_036\_TTLS \\
\midrule
LOC     & 1{,}883 & 8{,}868 \\
NUM     & 397     & 2{,}946 \\
PER     & 1{,}421 & 2{,}659 \\
TME     & 278     & 1{,}149 \\
TITLE   & 294     & 438 \\
DYNASTY & 171     & 415 \\
ORG     & 28      & 327 \\
\bottomrule
\end{tabular}
\caption{Phân bố nhãn thực thể theo từng quyển.}
\end{table}

\subsection{Cải tiến chất lượng NER}

Rà soát ngẫu nhiên các câu có nhãn LOC phát hiện ba lỗi hệ thống, đều đã được
sửa trong \texttt{pipeline/ner.py}:

\begin{itemize}
  \item \textbf{Chính tả gạch nối kiểu cũ không được nhận diện}: quyển LLKN
    dùng chính tả cũ nối các âm tiết bằng gạch nối (\texttt{Nguyễn-vương},
    \texttt{Gia-định}, \texttt{tri-phủ}), trong khi biểu thức chính quy nhận
    diện "từ viết hoa" chỉ chấp nhận một âm tiết liền, khiến các luật
    TME/DYNASTY/TITLE/LOC/PER dựa trên mẫu này bỏ sót toàn bộ tên riêng dạng
    này. Đã mở rộng để chấp nhận âm tiết nối gạch ngang.
  \item \textbf{Mô hình chung gán nhầm người có chức danh thành LOC/ORG}:
    xác nhận trực tiếp qua đầu ra của \texttt{underthesea.ner} (ví dụ
    \texttt{"Pagiơ"} sau \texttt{"Đô đốc"} bị gán \texttt{B-LOC}). Đã thêm một
    bước hậu xử lý: nếu một span LOC/ORG (hoặc phần văn bản ngay trước nó) bắt
    đầu/kết thúc bằng một chức danh đã biết trong từ điển, phần tên còn lại
    được gán lại nhãn PER.
  \item \textbf{Span vắt qua dấu câu giữa câu}: một số span của mô hình lấn
    qua dấu phẩy (ví dụ \texttt{"Tổng thống, Giôn-xơn"} bị gộp làm một span).
    Đã thêm bước cắt span tại dấu phẩy/chấm phẩy nằm giữa, chỉ giữ lại phần
    sau cùng.
\end{itemize}

Kết quả: số thực thể LOC giảm (LLKN: 2.049$\to$1.883; TTLS: 9.134$\to$8.868),
số thực thể PER tăng tương ứng (LLKN: 1.229$\to$1.421; TTLS:
2.361$\to$2.659), phản ánh đúng hướng sửa lỗi phân loại nhầm LOC/PER thay vì
chỉ thay đổi số lượng ngẫu nhiên. Các trường hợp mô hình chung gán nhầm nhãn
mà không có chức danh liền kề (ví dụ tên riêng đứng độc lập không theo sau
một chức danh trong từ điển) vẫn còn tồn tại và nằm ngoài phạm vi đợt sửa này
(xem mục Hạn chế).

\subsection{Kết quả hiệu đính OCR bằng Gemini}

\begin{table}[h]
\centering
\begin{tabular}{lrrrr}
\toprule
Quyển & Tổng đoạn & Đã sửa/cache & Còn lỗi (giữ nguyên) & Tỉ lệ hoàn thành \\
\midrule
HVQ\_036\_TTLS & 394 & 394 & 0 & 100.0\% \\
HVQ\_036\_LLKN & 153 & 153 & 0 & 100.0\% \\
\bottomrule
\end{tabular}
\caption{Tiến độ hiệu đính OCR theo từng quyển.}
\end{table}

Cả hai quyển đều đã được hiệu đính hoàn toàn (394/394 và 153/153 đoạn, 0 lỗi
API còn sót), đạt được qua nhiều lượt chạy lại tận dụng cơ chế cache (mỗi lượt
chỉ gọi lại API cho các đoạn từng thất bại ở lượt trước). Trong quá trình xử
lý, tầng miễn phí của Gemini API nhiều lần bị vượt ngưỡng quota
(\texttt{429 RESOURCE\_EXHAUSTED}), khiến một số đoạn phải giữ nguyên văn bản
gốc tạm thời (\emph{fallback\_error}) ở các lượt chạy trung gian; những đoạn
này sau đó được xử lý lại thành công ở các lượt tiếp theo (bao gồm cả việc đổi
sang API key khác khi cần) cho đến khi đạt 100\% cho cả hai quyển.

\subsection{Ví dụ một bản ghi NER}

\begin{verbatim}
{
  "sentence_id": "HVQ_036_TTLS_000001",
  "sentence": "Ngày 23, làm lễ mở đọc chiếu thư của nhà Minh ở điện Kính Thiên.",
  "entities": [
    { "text": "Ngày 23", "label": "TME" },
    { "text": "nhà Minh", "label": "DYNASTY" },
    { "text": "điện Kính Thiên", "label": "LOC" }
  ]
}
\end{verbatim}

\section{Cấu trúc thư mục output}

\begin{lstlisting}
output/
`-- HVQ_036/
     |-- HVQ_036_LLKN/
     |    |-- HVQ_036_LLKN_seg.tsv
     |    |-- HVQ_036_LLKN_ner.json
     |    |-- HVQ_036_LLKN_corrections.jsonl
     |    `-- cache/gemini.json
     `-- HVQ_036_TTLS/
          |-- HVQ_036_TTLS_seg.tsv
          |-- HVQ_036_TTLS_ner.json
          |-- HVQ_036_TTLS_corrections.jsonl
          `-- cache/gemini.json
\end{lstlisting}

\texttt{[matacpham]\_corrections.jsonl} không nằm trong đặc tả bắt buộc của đề
bài nhưng được giữ lại như một tài liệu phụ, ghi lại minh bạch từng đoạn văn
bản gốc/đã sửa và hành động tương ứng (\texttt{corrected}, \texttt{cache\_hit},
\texttt{fallback\_error}), phục vụ việc kiểm tra chất lượng hiệu đính.

\section{Hạn chế và hướng phát triển}

\begin{itemize}
  \item \textbf{Quota API miễn phí}: là nút thắt chính trong quá trình hiệu
    đính, đòi hỏi nhiều lượt chạy lại (tận dụng cache) và đổi API key khi cần
    để đạt 100\% cho cả hai quyển. Với tập dữ liệu lớn hơn, nên cân nhắc dùng
    gói trả phí để rút ngắn thời gian xử lý.
  \item \textbf{NER dựa một phần trên luật/gazetteer}: được tinh chỉnh riêng
    cho văn bản lịch sử Việt Nam trong tập dữ liệu này, có thể cần bổ sung
    thêm gazetteer khi áp dụng cho quyển sách khác ngoài phạm vi đồ án.
  \item \textbf{Mô hình CRF chung đôi khi gán nhầm nhãn không kèm chức danh}:
    ví dụ tên riêng đứng một mình (không theo sau chức danh trong từ điển)
    vẫn có thể bị gán nhầm LOC/ORG, hoặc mô hình tạo ra span dài bất thường
    gộp cả cụm danh từ lẫn động từ. Hướng khắc phục: dùng thông tin từ loại
    (POS) do underthesea cung cấp để kiểm tra tính hợp lệ khi mở rộng span
    theo nhãn BIO, hoặc thay bằng một mô hình NER được huấn luyện riêng cho
    văn bản lịch sử (ngoài phạm vi đợt cải tiến này).
  \item \textbf{Tách câu}: phụ thuộc vào chất lượng chuẩn hóa văn bản ở bước
    trước; các câu bị OCR làm hỏng nặng dấu câu có thể vẫn bị tách sai dù đã
    qua hiệu đính LLM.
\end{itemize}

\section{Kết luận}

Pipeline đã đáp ứng đầy đủ định dạng output theo yêu cầu đề bài (\texttt{seg.tsv},
\texttt{ner.json} với bộ nhãn tối thiểu PER/LOC/ORG/TITLE/TME/NUM và nhãn mở
rộng DYNASTY) cho cả hai quyển thuộc mã tác phẩm \texttt{HVQ\_036}, tổ chức
theo đúng cấu trúc thư mục nhiều quyển của đề bài. Cả hai quyển LLKN và TTLS
đều đã hoàn thành 100\% hiệu đính OCR bằng Gemini.

\end{document}
